\section{Summary} \label{sec:dykstra_summary}

Dykstra's algorithm provides a computationally lightweight method for calculating Euclidean projections, yet the unpredictable stalling phenomenon has historically inhibited its deployment in time-critical and high-frequency engineering frameworks~\cite{DYKSTRASTALLING}. This chapter successfully resolved the stalling problem for polyhedral feasible regions. By formalising the cyclical mechanics of the primary and auxiliary variables during a stall (Section~\ref{sec:stalling_phenomenon}), we derived an exact, closed-form mathematical solution for its duration (Section~\ref{sec:closed_form_solution}). This derivation directly enabled the construction of a fast-forward modified solver (Section~\ref{sec:fast_forward_algorithm}) that actively monitors the stationarity condition and instantly advances the internal memory variables past the stalling cycles.

As demonstrated by the empirical benchmarks (Section~\ref{sec:dykstra_numerical_experiments}), this algorithmic modification eliminates computational gridlock across diverse pathological geometries. It systematically reduces wall-clock execution time and accelerates convergence while strictly preserving the asymptotic convergence guarantees of the baseline method (Corollary~\ref{cor:convergence}). The theoretical derivations and optimisation advancements detailed throughout this chapter have been formally published as a preprint as preprint in~\cite{Vestini2025}.

With this core mathematical foundation, the fast-forward projection solver is now transformed into a highly predictable, real-time computational engine for computing Euclidean projections. The thesis now pivots from pure optimisation to the applied aerospace domain. The subsequent chapters will leverage this fast-forward architecture to enforce structural constraints within the PGD learning framework, enabling the accurate tracking of highly non-linear satellite uncertainty in the cislunar environment via optimal transport.